% ============================================================================
\section{Limitations and Conclusion}
\label{sec:conclusion}
% ============================================================================

The study covers one student scale, four tasks, and one system layout. GPAS
is a one-step rule: it treats micro-batches within a step as independent,
holds the AdamW scaling fixed within the step, and its variance reduction is
bounded by the count limits ($H\leq2$). The additive step-time model was
fitted to serial teacher execution; concurrent teacher serving needs its own
fit. The method needs at least two micro-batches per task per step, so
settings with many more tasks than micro-batches need a task-subsampling
variant.

GPAS fixes the objective weights and adapts only how many micro-batches each
task contributes, so the allocation changes gradient variance and nothing
else. The counts follow Neyman allocation in AdamW coordinates, Cost-GPAS adds
measured time, and both use statistics that MOPD training already produces.
\missing{Replace this sentence with the main four-teacher empirical
conclusion.}
